% concatenated from just-finite-11May2026.tex
%\documentclass[oneside]{amsart}
\documentclass[11pt,oneside]{amsart}
\usepackage{amsmath, amssymb, amsthm}

\newtheorem{theorem}{Theorem}
\newtheorem{lemma}[theorem]{Lemma}
\newtheorem{definition}[theorem]{Definition}
\newtheorem{proposition}[theorem]{Proposition}

\title{Every Finite Group Admits a Just Finite Presentation}
\author{Marc Lackenby}
%\date{\today}
\address{Mathematical Institute, University of Oxford, \newline Woodstock Road, Oxford OX2 6GG, United Kingdom}
\thanks{The author was partially supported by EPSRC grant EP/Y004256/1. For the purpose of open access, the author has applied a CC BY public copyright licence to any author accepted manuscript arising from this submission. Significant assistance was provided by the Google DeepMind AI co-mathematician tool. See the Methodology section for more details.}

\begin{document}

\maketitle

\begin{abstract}
A finite presentation $\langle X \, | \, R \rangle$ of a finite group $G$ is called ``just finite'' if removing any relation from $R$ results in a presentation for an infinite group. It has been an open question (Kourovka Notebook, Problem 21.10) whether every finite group admits such a presentation. We resolve this conjecture in the affirmative. 
%The proof proceeds by splitting into two cases: cyclic groups, which trivially admit balanced just finite presentations, and non-cyclic groups, for which we provide an explicit construction utilizing B.H. Neumann's presentation of the trivial group. 
\end{abstract}

\section{Introduction}
Inspired by the extensive theory of just infinite groups \cite{McCarthy1, McCarthy2, Wilson, Wilson2000, Grigorchuk2000, Reid:Structure, Reid2010}, Barnea introduced a notion that is, in some sense, dual \cite{Barnea}. He defined a \emph{just finite} presentation of a finite group $G$, which is a presentation $\langle X \, | \, R \rangle$ for the group with the property that, for every $r \in R$, the group $\langle X \, | \, R \setminus \{r\} \rangle$ is infinite. In  Problem 21.10 of the Kourovka Notebook \cite{Kourovka}, Barnea asked whether every finite group admits a just finite presentation. The goal of this note to prove that this is indeed the case.

\begin{theorem}
\label{Thm:Main}
Every finite group admits a finite, just finite presentation.
\end{theorem}

This is actually a consequence of a stronger result.
Recall that a group $G$ has \emph{Property (FA)} if any action of $G$ on a tree has a fixed point. It is a fundamental theorem of Serre \cite{Serre} that a finitely generated group has
Property (FA) if and only if it is neither a non-trivial amalgamated free product nor an HNN extension. We say that a presentation $\langle X \, | \, R \rangle$ for a group with Property (FA) is \emph{just-(FA)} if, for every $r \in R$, the group $\langle X \, | \, R \setminus \{r\} \rangle$ does not have Property (FA).

\begin{theorem}
\label{Thm:JustFA}
Every finitely presented group with Property (FA) admits a finite just-(FA) presentation.
\end{theorem}

This implies Theorem \ref{Thm:Main}, since a finite group $G$ clearly has Property (FA). Hence, it admits a finite just-(FA) presentation, which is therefore just-finite.

We also note the following consequence, concerning groups $G$ with Kazhdan's Property (T) \cite{Kazhdan}. We say that a presentation $\langle X \, | \, R \rangle$ for such a group is \emph{just-(T)} if, for every $r \in R$, the group $\langle X \, | \, R \setminus \{r\} \rangle$ does not have Property (T). Since Property (T) implies Property (FA) for countable discrete groups \cite{Watatani}, Theorem \ref{Thm:JustFA} gives the following result.

\begin{theorem}
\label{Thm:JustT}
Every finitely presented group with Property (T) admits a finite just-(T) presentation.
\end{theorem}


\section{The Construction}

The aim is to start with a finite presentation $\langle X \, | \, R \rangle$ of the given group $G$ with Property (FA) and modify it to a just-(FA) presentation. Each relation will be substituted for two new relations that have a form based on the following lemma.

\begin{lemma}[B.H. Neumann \cite{Neumann}]
\label{lem:neumann}
Let $H$ be a group with elements $u, v \in H$ satisfying $u^{-1}vu = v^2$ and $v^{-1}uv = u^2$. Then $u=1$ and $v=1$.
\end{lemma}

\begin{proof}
From the second relation, $v^{-1}uv = u^2$, we can write the commutator $[u, v] = u^{-1}(v^{-1}uv) =
u^{-1}u^2 = u$. From the first relation, $u^{-1}vu = v^2$, we can invert it to get $u^{-1}v^{-1}u =
v^{-2}$. We can then write the commutator as $[u, v] = (u^{-1}v^{-1}u)v = v^{-2}v = v^{-1}$. Equating
the two expressions for the commutator, we have $u = v^{-1}$. Substituting this into the first relation
$u^{-1}vu = v^2$ yields $(v^{-1})^{-1}v(v^{-1}) = v^2$, which simplifies to $v = v^2$. This implies $v = 1$, and therefore
$u = 1$.
\end{proof}

The following type of semi-direct product of cyclic groups plays a technical role in our proof.

\begin{lemma}
\label{Lem:SemiDirect}
For $k > 1$, the presentation $\langle x, y \, | \, x^{-1} y x = y^2, x^k = 1 \rangle$ specifies the group $\mathbb{Z}_{2^k-1} \rtimes \mathbb{Z}_k$, where $\langle x \rangle = \mathbb{Z}_k$ and $\langle y \rangle = \mathbb{Z}_{2^k-1}$ and the conjugation action of $x$ on $\langle y \rangle$ is $y \mapsto y^2$.
\end{lemma}

\begin{proof}
There is a homomorphism $\langle x, y \, | \, x^{-1} y x = y^2, x^k = 1 \rangle \rightarrow \mathbb{Z}_{2^k-1} \rtimes \mathbb{Z}_k$, since the two relations are satisfied in the image group. It is surjective because the images of $x$ and $y$ generate $\mathbb{Z}_{2^k-1} \rtimes \mathbb{Z}_k$. Any element in $\langle x, y \, | \, x^{-1} y x = y^2, x^k = 1 \rangle$ can be written as $x^m y^n$ for some integers $m$ and $n$. Using $x^k = 1$, we may assume $0 \leq m < k$. Conjugating $y$ by $x$ a total of $k$ times gives $y^{2^k}$. Since $x^k = 1$, we deduce that $y^{2^k} = y$. So, we may also assume $0 \leq n < 2^{k}-1$. So $\langle x, y \, | \, x^{-1} y x = y^2, x^k = 1 \rangle$ has at most $k(2^k -1)$ elements. However, $\mathbb{Z}_{2^k-1} \rtimes \mathbb{Z}_k$ has exactly this many elements, and therefore the above surjective homomorphism is an isomorphism.
\end{proof}

We are now in a position to prove our main theorem.

\begin{proof}[Proof of Theorem \ref{Thm:JustFA}]
Let $G$ be a finitely presented group with Property (FA). If $G$ is cyclic, $G$ must be a finite cyclic group $\mathbb{Z}_n$ and so admits the balanced, just-(FA) presentation $\langle x \, | \, x^n = 1 \rangle$.

Assume $G$ is non-cyclic. Let $\langle X \, | \, R \rangle$ be an irredundant finite presentation for $G$ (that is, for any $r \in R$, the group $H_r = \langle X \, | \, R \setminus \{r\} \rangle$ is not isomorphic to $G$, implying $r \neq 1$ in $H_r$). 

We define a new presentation $P' = \langle X' \, | \, R' \rangle$ as follows:
\begin{itemize}
    \item $X' = X \cup \{b_r \, | \, r \in R\}$
    \item $R' = \{r^{-1} b_r r = b_r^2 \, | \, r \in R\} \cup \{b_r^{-1} r b_r = r^2 \, | \, r \in R\}$
\end{itemize}
By Lemma \ref{lem:neumann}, these relations force $b_r = 1$ and $r = 1$ for all $r \in R$. (See Section \ref{Sec:Examples} for a specific example of such a presentation $P'$.)

We claim that $P'$ presents $G$. Consider what happens to the group $G = \langle X \, | \, R \rangle$ when a new generator $b_r$ is added, the relation $r$ is removed and then replaced by 
the two new relations $r^{-1} b_r r = b_r^2$ and $b_r^{-1} r b_r = r^2$. These force $b_r = 1$ and $r = 1$. Conversely $b_r = 1$ and $r = 1$ imply the two new relations 
$r^{-1} b_r r = b_r^2$ and $b_r^{-1} r b_r = r^2$. Thus, instead of adding these two new relations, we could have simply set $b_r = 1$ and $r = 1$. In other words, the new group is obtained from $G = \langle X \, | \, R \rangle$ by adding the new generator $b_r$ and a new relation $b_r = 1$. Hence, this does not change the group. Repeating this for each relation in $R$, we deduce that $P'$ does present $G$.

We must show $P'$ is just-(FA). Let $r \in R$. We show that removing either of the corresponding relations yields a group without Property (FA).

\textbf{Case 1: Remove $b_{r}^{-1} r b_{r} = r^2$.}
Let $K_1$ be the resulting group. As above, let $H_r$ be the group $\langle X \, | \, R \setminus \{ r \} \rangle$. This also has a presentation obtained from $P'$ by removing the generator $b_r$ and the relations  $r^{-1} b_{r} r = b_{r}^2$ and $b_r^{-1} r b_r = r^2$. Let $k \in \mathbb{N} \cup \{ \infty \}$ be the order of $r$ in $H_r$. Since $R$ is irredundant, $k > 1$. Thus $K_1$ can be decomposed as an amalgamated free product:
\[ K_1 \cong H_{r} *_{\langle r \rangle = \langle x \rangle} B' \]
where $B' = \langle x, b_{r} \, | \, x^{-1} b_{r} x = b_{r}^2, x^k = 1 \rangle$. When $k = \infty$, we do not include the relation $x^k = 1$. Thus, when $k = \infty$, $B'$ is the Baumslag-Solitar group $B(1,2)$. When $k$ is finite,
Lemma \ref{Lem:SemiDirect} gives that the group $B'$ is the semi-direct product $\mathbb{Z}_{2^k-1} \rtimes \mathbb{Z}_k$, and the index of the amalgamating subgroup $\langle x \rangle$ in $B'$ is $2^k - 1$. Thus, when $k = \infty$ or $k \ge 2$, this index is at least $3$.

An amalgamated free product $A *_C B$ is non-trivial when $C \not= A$ and $C \not= B$.
Since $[B' : \langle x \rangle] \ge 3 > 1$, $K_1$ is a non-trivial amalgamated free product unless $H_{r} = \langle r \rangle$.
In this case, $H_{r}$ is a cyclic group. However, $G$ is the quotient of $H_{r}$ by the normal closure of $r$. Since any quotient of a cyclic group is cyclic, this contradicts the assumption that $G$ is non-cyclic. Therefore, $H_{r} \neq \langle r \rangle$, and $K_1$ is a non-trivial amalgamated free product.

\textbf{Case 2: Remove $r^{-1} b_{r} r = b_{r}^2$.}
Let $K_2$ be the resulting group. The remaining relation is $b_{r}^{-1} r b_{r} = r^2$.
We define a homomorphism $\phi: K_2 \to \mathbb{Z} = \langle t \rangle$ by setting $\phi(x)=1$ for all $x \in X$, $\phi(b_{r'})=1$ for $r' \neq r$, and $\phi(b_{r})=t$. Note that since $X$ maps to the identity, all original relations in $R \setminus \{r\}$ are trivially satisfied. Under this map, the remaining relation maps to $t^{-1} \cdot 1 \cdot t = 1$. This homomorphism is well-defined and maps onto the infinite group $\mathbb{Z}$. Hence, $K_2$ does not have Property (FA).

Thus, removing any relation from $P'$ yields a group without Property (FA), and so $P'$ is a just-(FA) presentation for $G$.
\end{proof}

We remark, as a consequence of the above proof, the following result. Recall that the \emph{deficiency} of a presentation $\langle X \, | \, R \rangle$ is $|X| - |R|$.

\begin{theorem}
Let $P$ be any finite irredundant presentation of a non-cyclic group $G$ with Property (FA). Then $G$ admits a finite, just-(FA) presentation $P'$ where the deficiency of $P'$ is equal to that of $P$, and where the generators of $P'$ represent the same elements of $G$ that the generators of $P$ did, plus multiple copies of the identity element.
\end{theorem}

\section{Examples}
\label{Sec:Examples}
The following is a presentation of the dihedral group of order 8:
$$\langle \sigma, \tau \, | \, \sigma^4 = 1, \tau^2 = 1, \tau^{-1} \sigma \tau = \sigma^3 \rangle.$$
When the first relation is removed, this gives a semi-direct product $\mathbb{Z}_{8} \rtimes \mathbb{Z}_2$, by an argument analogous to the one in Lemma \ref{Lem:SemiDirect}. When the second relation is removed, the resulting group admits a surjective homomorphism onto $\mathbb{Z} = \langle t \rangle$, sending $\sigma$ to $1$ and $\tau$ to $t$. When the third relation is removed, we have the free product $\mathbb{Z}_4 \ast \mathbb{Z}_2$. So this presentation is irredundant but not just finite. When we apply the above procedure to it, we obtain the just finite presentation
$$\left \langle \sigma, \tau, a, b, c \;\middle|\;
\begin{array}{c} \sigma^{-4} a \sigma^4 = a^2 , a^{-1} \sigma^4 a = \sigma^8, \\
\tau^{-2} b \tau^2 = b^2 , b^{-1} \tau^2 b = \tau^4, \\
(\tau^{-1} \sigma \tau \sigma^{-3})^{-1} c (\tau^{-1} \sigma \tau \sigma^{-3}) = c^2,\\
c^{-1} (\tau^{-1} \sigma \tau \sigma^{-3}) c= (\tau^{-1} \sigma \tau \sigma^{-3})^2 
\end{array}
\right \rangle.$$
However, the following presentation for the dihedral group of order 8 is already just finite:
$$\langle \sigma, \tau \, | \, \sigma^4 = 1, \tau^2 = 1, \tau^{-1} \sigma \tau = \sigma^{-1} \rangle.$$
Removing the first relation gives the infinite dihedral group.

\section{Methodology}
This paper was produced using the AI co-mathematician tool \cite{Zheng} developed by Google DeepMind. The author posed Problem 21.10 from the Kourovka Notebook to the AI co-mathematician. AI co-mathematician proposed a solution, which was essentially the proof given in this paper, but then found a flaw in its argument which it was unable to resolve. The issue was that it was uncertain how to proceed in the case where $H_r = \langle r \rangle$  in Case 1 of the proof of Theorem \ref{Thm:JustFA}. By analysing the argument carefully, the author was able to find a fix for it, which the AI co-mathematician then confirmed was correct and wrote up. The author then made some minor edits and clarifications. The author then introduced the notions of just-(FA) and just-(T) presentations, and made the necessary minor adjustments to the proofs, resulting in this paper. The paper was then checked and corrected by the AI co-mathematician tool.

\bibliography{just-finite-biblio}
\bibliographystyle{plain}


\end{document}
